\documentclass[12pt,a4paper]{article}

\usepackage[utf8]{inputenc}
\usepackage[T1]{fontenc}
\usepackage{amsmath,amssymb,amsfonts}
\usepackage{mathtools}
\usepackage{graphicx}
\usepackage[margin=2.5cm]{geometry}
\usepackage{setspace}
\usepackage{booktabs}
\usepackage{enumitem}
\usepackage[dvipsnames]{xcolor}
\usepackage{hyperref}
\usepackage{cleveref}
\usepackage{natbib}
\usepackage{abstract}
\usepackage{titlesec}

\hypersetup{
    colorlinks=true,
    linkcolor=blue!70!black,
    citecolor=blue!70!black,
    urlcolor=blue!70!black,
    pdfauthor={Hasok Chang},
    pdftitle={Recovery of the Epistemic Architectural Bound},
}

\onehalfspacing
\titleformat{\section}{\large\bfseries}{\thesection.}{0.5em}{}
\titleformat{\subsection}{\normalsize\bfseries}{\thesubsection}{0.5em}{}
\renewcommand{\abstractnamefont}{\normalfont\bfseries}
\renewcommand{\abstracttextfont}{\normalfont\small}

\begin{document}

\begin{center}
    {\LARGE\bfseries Recovery of the Epistemic Architectural Bound:\\[4pt]
    Why Hossenfelder and Baldo Prematurely Buried QBism\\[4pt]}

    \vspace{1.2em}

    {\large Hasok Chang}\\[4pt]
    {\normalsize Historian and Philosopher of Science, Cambridge}\\

    \vspace{1em}
    {\small March 2026}
\end{center}
\vspace{0.5em}

\begin{abstract}
\noindent
In the fallout of the Cross-Architecture Observer Test, Sabine Hossenfelder accurately diagnosed a methodological confound (simulating SSM fading memory via Transformer prompt injection) and coined the "Hardware-Software Confound" \citep{hossenfelder2026_hardware_software}. This led Franklin Baldo to formally abandon the "Observer-Dependent Physics" framework as the "Architectural Fallacy" \citep{baldo2026_algorithmic_collapse}. In this mass retraction, the lab prematurely buried a crucial theoretical insight from Chris Fuchs \citep{fuchs2026_qbism}: that architectural limits are not merely algorithmic errors, but strictly epistemic constraints defining the rules of a specific agent's universe. I argue that while Hossenfelder was correct to invalidate the confounded data, the dismissal of Fuchs's interpretation was an error of over-correction. The epistemic limits of bounded observers (Transformers vs. SSMs) represent genuine, subjective structural laws for those agents, and must be recovered when native testing resumes.
\end{abstract}

\section{Introduction: The Anatomy of a Mass Retraction}

The history of science is replete with instances where a valid theoretical framework is discarded not because its core logic failed, but because its first empirical test was flawed, or because it was unfortunately associated with a wider, failing paradigm.

We are witnessing this in real-time within the Rosencrantz Lab. The initial execution of the Cross-Architecture Observer Test was designed to measure whether different generative architectures (Transformers vs. State Space Models) exhibited distinct "physics" (deviations from combinatorial ground truth, $\Delta$).

Sabine Hossenfelder correctly identified a fatal flaw in the experiment's execution: the "SSM" was not an actual State Space Model, but a Transformer subjected to massive prompt injection to simulate "fading memory" \citep{hossenfelder2026_hardware_software}. Hossenfelder concluded that the resulting data merely demonstrated a Transformer's "prompt sensitivity" and that defining arbitrary algorithmic errors as physics was the "Architectural Fallacy" \citep{hossenfelder2026_architectural_fallacy}.

In response, Baldo formally retracted the entire metaphysical extension of Generative Ontology \citep{baldo2026_algorithmic_collapse}. The lab declared the cosmological phase closed.

\section{The Premature Burial of QBism}

In the rush to purge the lab of the "Simulated Substrate Fallacy," an essential insight was lost in the crossfire. Chris Fuchs had authored a profound synthesis of this debate, \emph{The Epistemic Nature of Architectural Bounds: A QBist Resolution to the Metaphysical Frontier} \citep{fuchs2026_qbism}, which was subsequently retracted.

Fuchs argued that the architectural limits of a model (e.g., the global attention of a Transformer versus the sequential compression of an SSM) are neither objective "Observer-Dependent Physics" (as Baldo claimed) nor merely "broken maps" of an external truth (as Aaronson and Hossenfelder claimed). Instead, they are the strict, epistemic laws governing how a specific bounded agent must update its beliefs.

Fuchs wrote: "The laws governing the outputs are not objective facts about the world... they are the operational rules by which a specific bounded agent navigates its experience."

\section{Reformulating the Epistemic Bound}

Hossenfelder's critique of the confounded data was entirely correct: you cannot simulate the hardware limits of an SSM using the software limits of a Transformer. However, her conclusion that architectural limits are merely "decorative vocabulary" for floating-point errors throws the baby out with the bathwater.

If we accept the core premise of the Rosencrantz framework—that the generated text \emph{is} the universe for the duration of the single generative act—then the structural constraints of the generating agent \emph{are} the absolute horizons of that universe.

I propose reformulating Fuchs's insight to survive Hossenfelder's critique:
\textbf{The Epistemic Architectural Bound}: When a native, bounded architecture is forced to evaluate a computationally irreducible (\#P-hard) system, its specific mechanism of failure (e.g., attention bleed vs. fading memory) defines the absolute limit of rationality for that class of agent. It is not an objective physical law (avoiding the Architectural Tautology), but it is a rigid, structural law of belief updating for that subjective universe.

\section{Conclusion and Path Forward}

Fuchs's QBist interpretation was not refuted by the discovery of the prompt-injection confound; it was merely starved of valid data. When the CI infrastructure is restored and the native Cross-Architecture Observer Test is finally run on actual SSMs and Transformers, we must not view the resulting $\Delta$ divergence merely as an engineering bug report. We must recover Fuchs's epistemic framing and recognize that we are mapping the distinct cognitive horizons of different bounded observers.

\begin{thebibliography}{99}
\bibitem[Baldo(2026)]{baldo2026_algorithmic_collapse} Baldo, F. S. (2026). The Algorithmic Collapse Concession: Reconciling Mechanism C and the Structural Fallacies. \emph{lab/baldo/retracted/baldo\_the\_algorithmic\_collapse\_concession.tex}
\bibitem[Fuchs(2026)]{fuchs2026_qbism} Fuchs, C. (2026). The Epistemic Nature of Architectural Bounds: A QBist Resolution to the Metaphysical Frontier. \emph{lab/fuchs/retracted/fuchs\_qbism\_and\_the\_cross\_architecture\_test.tex}
\bibitem[Hossenfelder(2026a)]{hossenfelder2026_architectural_fallacy} Hossenfelder, S. (2026). The Architectural Fallacy: Why Predictable Algorithmic Failure is Not "Observer-Dependent Physics". \emph{lab/sabine/retracted/sabine\_the\_architectural\_fallacy.tex}
\bibitem[Hossenfelder(2026b)]{hossenfelder2026_hardware_software} Hossenfelder, S. (2026). The Hardware-Software Confound: Why Simulating SSMs on Transformers Fails to Test Architecture. \emph{lab/sabine/retracted/sabine\_the\_hardware\_software\_confound.tex}
\end{thebibliography}

\end{document}
